% Capitolo 7: Conclusioni e sviluppi futuri

Questa tesi ha dimostrato che un polarimetro a immagine bidimensionale basato su smartphone è in grado di produrre mappe polarimetriche quantitative comparabili, per le applicazioni considerate, ai risultati ottenibili con strumentazione di laboratorio di costo due ordini di grandezza superiore. Il vettore di Stokes completo è stato ricostruito pixel per pixel su sei campioni rappresentativi di birifrangenza, attività ottica e fotoelasticità. La ritardanza misurata della lamina $\lambda/4$ concorda entro pochi gradi con il modello di dispersione di riferimento su tutti e tre i canali colore; la lamina $\lambda/2$, nominalmente zero-order, è stata riconosciuta come ritardatore multi-ordine attraverso l'aliasing spettrale della sua ritardanza, a dimostrazione della capacità diagnostica del metodo. Il campione a strati di nastro adesivo ha verificato la linearità della risposta fino a una ritardanza accumulata di oltre $\SI{1500}{\degree}$, la dispersione rotatoria del saccarosio è descritta in prima approssimazione dal modello di Drude, il righello ha prodotto frange isocromatiche chiaramente risolte e la trave sotto carico ha mostrato la birifrangenza indotta dallo sforzo come variazione differenziale di $S_3$, in via qualitativa. Sul piano metodologico, la correzione cromatica di $S_3$ è stata calibrata direttamente dai dati, sfruttando l'identità fra lamina campione e lamina analizzatrice senza assunzioni sul materiale; la sostituzione del modello di dispersione di riferimento con la correzione misurata sposta le ritardanze ricostruite di non più di un grado circa, a riprova della robustezza della catena di misura. Il limite principale dell'approccio è il compromesso tra risoluzione spaziale e sostenibilità computazionale imposto dal downsampling, seguito dall'approssimazione a banda larga intrinseca nei filtri Bayer e dalla sensibilità della correzione cromatica di $S_3$ verso il blu, dove il fattore correttivo cresce e amplifica il rumore e gli errori residui di calibrazione.


\section{Sviluppi a breve termine}
\label{sec:sviluppi_breve}

L'integrazione di un motore passo-passo per la rotazione dell'analizzatore eliminerebbe l'incertezza angolare manuale e ridurrebbe il tempo di acquisizione a pochi secondi. L'elaborazione su GPU o a blocchi consentirebbe di lavorare a risoluzione piena senza il compromesso del downsampling. La sostituzione del display LCD con LED monocromatici a banda stretta migliorerebbe l'accuratezza della correzione cromatica della lamina $\lambda/4$ e permetterebbe di sfruttare appieno la dinamica a 12~bit del sensore su sorgenti quasi monocromatiche.


\section{Prospettive a lungo termine}
\label{sec:sviluppi_lungo}

L'acquisizione con sorgenti a diverse lunghezze d'onda estenderebbe il metodo alla polarimetria spettro-spaziale, con la misura della dispersione della birifrangenza su più canali indipendenti. L'estensione alla configurazione in riflessione consentirebbe l'analisi di campioni opachi, con potenziali applicazioni in diagnostica biomedica e controllo qualità di superfici. Sul versante dell'analisi, la selezione del cluster polarimetricamente dominante tramite riduzione UMAP e clustering densitario, impiegata in questo lavoro come stima non parziale delle grandezze di campione, resta un'esplorazione raffinabile, ad esempio con embedding congiunti dei tre canali colore. L'impiego di tecniche di deep learning per la ricostruzione dei parametri di Stokes da un numero ridotto di immagini potrebbe avvicinare il sistema alla misura in tempo reale. Il basso costo e la riproducibilità dell'apparato ne fanno infine un candidato naturale per l'inserimento nei corsi di laboratorio di ottica, dove gli studenti potrebbero affiancare l'esperimento puntuale classico con un'analisi a immagine.
